% Abstract — faithful mirror of pretext/resumo.tex (NBR 6028: one
% paragraph, 150-500 words, no formulas).
For a decade, model choice for tabular data was guided by the maxim that
gradient boosting always wins; the arrival of tabular foundation models
--- pre-trained networks that predict with no task-specific training ---
made that guidance obsolete without replacing it with neutral evidence.
This dissertation attacks the problem on three chained fronts. First, a
multi-criteria benchmark evaluates 14 models from three families
(gradient boosting, deep learning, and foundation models) on 18 public
OpenML datasets (250 of the 252 possible cells; two declared structural
exclusions), under nested cross-validation with an untouched hold-out and
both frequentist and Bayesian statistical inference. In the study's 11-model core slate, the top of the accuracy ranking is a statistical tie --- 8 of
the 10 non-leading models are practically equivalent to the leader under
a region of practical equivalence of one percentage point of AUC --- and
the 2026 generation breaks it: TabFM, evaluated zero-shot, leads the
average rank in all three tasks (2.7 in binary classification, 1.0 in
multiclass, and 1.2 in regression) and is the best model on 13 of the 18
datasets, at the cost of four to five orders of magnitude in inference
latency; Kolmogorov--Arnold networks, independently tested, finish in the
bottom third across all three tasks. Second, a decision framework is
validated by leave-one-dataset-out: meta-feature routing does not
generalize (hit rate of 0.11 against 0.44 for the trivial baseline), the
fixed ``foundation model first'' policy is regret-optimal on the primary
metric (median regret of 0.000 in the 2026 generation), and a
pre-registered multi-criteria regret experiment shows that shifting just
10 to 20 percent of the weight from accuracy to latency (crossover
between weights 0.9 and 0.8) flips the recommendation in favor of
gradient boosting --- quantitatively justifying the flowchart of four
a priori verifiable constraints that consolidates the framework. Third,
the distributional distillation of foundation models --- a scope verified
as novel --- transfers the quantile curves of the TabPFN teacher to
students that serve in microseconds: point distillation fails robustly at
realistic data scales, but distributional distillation retains a median
of 19 percent of the teacher's distributional advantage, positive on 16
of 20 datasets (sign test with a p-value of 0.006), beats classical
distributional baselines on 17 of 20 datasets in the pre-registered
experiment E1, and condenses into a practical decision rule whose
preconditions can be verified in minutes. Together --- evidence under a
single protocol, a validated framework, and selective distillation ---
these results replace the maxim with a procedure; and what the foundation
model knows that raw data cannot teach the student is the predictive
distribution --- that is what gets distilled.

\vspace{1cm}
\noindent\textbf{Keywords:} tabular data; foundation models;
benchmarking; knowledge distillation; distributional regression.
